% !PRINTED 619
\weberpage{619}
%% [page furniture, not body text -- source-confirmed 2026-08-10]
%% \S\,169. \quad Relation of the division fields to the class field. \quad 619

%<<CONT
and the relative degree with respect to $\Omega$ at most equal to $\tfrac{1}{2} h''\varphi(m)$, and from this it follows by (11) that $h''$ cannot be smaller than $H'_g$. Thus it is proved:

5. \gesperrt{The class equation $H_{-D}(u)$ is, in the field that contains all roots of unity, not further decomposable than into the factors corresponding to the genera.}

\noindent\textbf{\S\,169.\qquad \gesperrt{Relation of the division fields to the class field.}}

In the last two sections we have considered two kinds of groups, of which the first led to the division fields of the elliptic functions, the second to the order fields in combination with roots of unity. Between these there exists a relation:

The first group $A$ consisted of the numbers $\alpha$ in $O$ which satisfied the condition
\weq{1}{\alpha \equiv \pm 1 \pmod{\mathfrak{m}}}
where $\mathfrak{m}$ is any ideal in $\Omega$; the second group $A'$ consisted of the numbers $\alpha'$ of the order $[Q]$ which satisfied the condition
\weq{2}{N(\alpha') \equiv 1 \pmod{m}}
where $m$ and $Q$ denote natural numbers.

If we assume that $\mathfrak{m}$ is a prime ideal or a power of a prime ideal, then we can choose $Q$ and $m$ in such a way that the group $A'$ is contained in the group $A$.

For if $r$ is a rational number, then, since $\alpha'$ belongs to the order $[Q]$:
\weq{3}{\alpha' \equiv r \pmod{Q}, \quad N(\alpha') \equiv r^2 \pmod{Q},}
and if $Q$ is divisible by $m$:
\weq{4}{\alpha' \equiv r \pmod{m},}
hence by (2) and (3)
\[
r^2 \equiv 1 \pmod{m}
\]
\weq{5}{(r - 1)(r + 1) \equiv 0 \pmod{m},}
we therefore assume $Q$ to be divisible by $m$, and $m$ by $\mathfrak{m}$; then by (5), if $\mathfrak{m}$, as we have assumed, is a prime ideal or a power of such a one, either
\weq{6}{r \equiv 1 \text{ or } r \equiv -1 \pmod{\mathfrak{m}},}
%CONT>>
